\documentclass[12pt]{article}

\usepackage[a4paper,margin=1in]{geometry}
\usepackage{times}
\usepackage[T1]{fontenc}
\usepackage[utf8]{inputenc}
\usepackage{enumitem}
\usepackage{titlesec}
\usepackage{xcolor}
\usepackage{hyperref}
\usepackage{setspace}

\onehalfspacing

\titleformat{\section}{\normalfont\Large\bfseries\color{blue!40!black}}{\thesection}{1em}{}
\titleformat{\subsection}{\normalfont\large\bfseries\color{blue!40!black}}{\thesubsection}{1em}{}

\begin{document}

\section*{Chapter 1 --- Introduction}
\addcontentsline{toc}{section}{Chapter 1 --- Introduction}

\subsection*{1.1. Background}

In recent years, IT infrastructure has increasingly been deployed in distributed, small- to medium-scale models --- typified by \textbf{Mini Data Centers} located at branch offices, local server rooms, or edge computing sites. Unlike large-scale Data Centers, which benefit from dedicated operations teams and highly automated monitoring systems, Mini Data Centers often lack sufficient personnel and tooling, making infrastructure monitoring and maintenance considerably more difficult.

At present, infrastructure monitoring relies mainly on \textbf{software-based monitoring systems} (dashboards displaying CPU, RAM, Disk, and server status). However, these systems reveal a clear gap between \textbf{physical infrastructure} (servers, racks, cabling, etc.) and \textbf{digital monitoring} (the digital data shown on screens): when an alert appears on the dashboard, a technician must still manually locate the corresponding physical device on-site --- typically by reading labels, tracing rack positions, or relying on personal experience.

The advancement of \textbf{Augmented Reality (AR)} technology opens up the possibility of bridging this gap. By scanning a physical device (via a QR/ArUco marker), a system can instantly identify that device and display real-time telemetry, alerts, and related tickets directly on a mobile screen. This is the underlying motivation for integrating AR into the infrastructure monitoring and maintenance process, rather than keeping the physical and digital worlds separate as they are today.

\subsection*{1.2. Problem Statement}

Building on this background, several concrete problems can be identified --- problems that the proposed system, \textbf{AR-IMMS (AR-Integrated Monitoring \& Maintenance System)}, aims to address:

\begin{itemize}[leftmargin=*]
    \item \textbf{Server data is scattered across the monitoring system}: Metrics such as CPU, RAM, Disk, and network status are collected and displayed on a dashboard, but they are not directly linked to the physical location of the corresponding device within a rack/room.

    \item \textbf{Technicians must manually identify the physical device on-site}: Upon receiving an alert or ticket, a technician must manually look up the correct device (based on Server ID, rack diagrams, or physical labels) --- a process prone to error, especially in environments with many similar-looking servers.

    \item \textbf{Alerts and tickets are not directly tied to device location}: When an alert is triggered, it is typically presented only as text/numeric data (Server ID, CPU level, etc.), with no visual mechanism linking it to the actual physical location --- slowing down response and increasing the risk of mistakes.

    \item \textbf{Manual inspection is time-consuming and inconsistent}: The current inspection process depends heavily on individual experience and lacks visual support tools, resulting in a longer Mean Time to Repair (MTTR).

    \item \textbf{Lack of integration between telemetry, Digital Twin, and AR}: Real-time monitoring data (telemetry) is not currently mapped into a digital model (Digital Twin) that can be accessed on-site through an AR device --- this is the technological gap the project seeks to fill.
\end{itemize}

Based on these issues, the project poses the following central research question: \textit{``How can digital monitoring data be connected to physical devices in a visual, instantaneous way, enabling technicians to identify and resolve issues faster and more accurately?''} --- AR-IMMS is the proposed solution to this question (detailed further in Section~1.3).

\end{document}
